\subsection{带风险控制和复用成本的输出一致性策略}
后续策略可在有效的选择后置信界下，最小化校准与预计复用时间，
并认证逐查询的\emph{参考一致性}损失。输入为匹配执行路径的逐query不一致曲线
和固定参考，变量包括样本量及survivor预算。绝对GT目标还需要独立标签验证，
或参考与GT之间的额外假设。平稳性、逐查询抽样与共同分数／排序契约必须明确，
漂移触发器应使旧证书失效并重新认证。

最低实验应在留出池比较固定样本规则、秩曲线复用和成本策略，记录总耗时与尾部
一致性损失。旧四单元离线回放的序贯控制器已因额外搜索launch吞掉样本节省而失败，
单纯重复序贯测试不足以形成改进。新的问题是：复用中间信息与质量证书能否联合
改善\emph{总}成本。违反风险限制，或比相同质量的固定样本对照没有省时，即失败。

\subsection{考虑容量的残差放置}
大型索引中常驻残差占主导。可研究primary常驻、残差block在显式显存上限内缓存，
miss从主机内存取回。固定编码、路由、survivor集合与精排算术，隔离放置与表示质量。
最低实验改变显存上限、batch和查询局部性，测hit、传输字节、重叠、尾延迟和吞吐。
只有在可接受的总时间前沿上扩展容量才算有效；原本可常驻规模上更慢的cache不足以
支持方案。这是未实现的存储扩展，与已有残差下界pilot不同；降低$\alpha$也不实现它，
因为完整数组依然分配。
